\section{Discussion}

Across four validation stages one property of the map recurs. Immunometabolic load is the axis
that stands apart from general severity ($\phi_{Gd}\approx 0.10$ versus $0.39$/$0.42$ for
cognition and sleep), the axis that stays put over two years (test--retest ICC $0.91$), and the
axis that marks the worst-prognosis pole of the patient space (two-year functional remission
$22\%$ at the immunometabolic corner versus $63\%$ at the well pole). No single one of these
observations would carry the argument; their convergence---separable \emph{and} durable \emph{and}
prognostic, the same axis each time---is what identifies the immunometabolic dimension as a
distinct, stable, outcome-relevant feature of severe mental illness rather than another face of
being severely ill.

This headline is corroborated, in direction, by a large and independent literature, which is why
we state it as \emph{relative} rather than absolute independence. Metabolic burden in severe
mental illness is driven by adiposity, medication, age and chronicity more than by diagnosis or
symptom severity, is present transdiagnostically across schizophrenia, bipolar disorder and
depression~\citep{vancampfort2015mets,vancampfort2016diabetes,pillinger2017,perry2019,garridotorres2022};
inflammation marks a subgroup and specific symptoms rather than the severity
continuum~\citep{osimo2019,fernandes2016bd,fernandes2016sz}; and the immunometabolic-depression
programme has independently arrived at a separable biological dimension mapping onto atypical,
energy-related symptoms~\citep{penninx2024imd,milaneschi2020,lamers2013,lamers2020,milaneschi2017}.
Our contribution is not to re-discover this biology but to \emph{measure} it inside one
diagnosis-blind map: to give the qualitatively-known separation a number ($\approx 1\%$ shared
variance with general burden), estimated with a model honest enough---no imputation, uncertainty
carried end-to-end, the map frozen before validation---that the number can be trusted.

The practical reading is a design principle, not a treatment rule. The axes divide cleanly into
those worth \emph{stratifying} on and those worth \emph{monitoring}: a durable, severity-%
independent axis that concentrates poor outcomes is what one enriches a trial or a service cohort
on, whereas the mobile symptom axes are what one tracks over time. The immunometabolic axis sits
squarely on the stratify side---it is still there at two years, it is not captured by a severity
score, and it identifies a corner of the map where only one patient in five reaches functional
remission. That makes it a concrete enrichment and monitoring target: a rationale for measuring
and following routine cardiometabolic and inflammatory biology as an axis of illness in its own
right, and for testing---prospectively---whether intervening on it changes course.

\medskip\noindent\textbf{Limitations.} Three classes bound these claims.
\emph{First, the results are observational and internal.} Every finding is an internal-validity
result on baseline and two-year data from a single consortium; the map's structure, separability
and prognostic gradient are estimated, not externally validated, and causal language is
deliberately avoided.
\emph{Second, the prognostic signal is group-level.} The archetype encoding improves held-out
predictive density for two-year functioning ($\Delta\text{ELPD}=+62.8$) but adds little
individual-level discrimination (AUC $+0.010$) and is co-informative with, not superior to,
diagnosis. The immunometabolic pole is an enrichment target, not an individual predictor, and we
report the modest effect as plainly as the strong ones.
\emph{Third, coverage and horizon are finite.} Follow-up outcomes were available for $2{,}420$ of
$9{,}013$ patients, the biology panel is the routine one collected across these cohorts rather than
a deep immunophenotype, and two years is a short window for functional trajectories.

Establishing individual-level clinical utility would require external validation in independent
cohorts, incident-event outcomes over longer horizons, and---for any causal claim about the
immunometabolic axis---prospective or interventional data beyond this baseline observational
cohort. What the present map establishes is prior to all of that and necessary for it: that when
biology is measured honestly and placed inside the transdiagnostic space as a co-equal axis, one
biological dimension emerges as distinct, durable and prognostic---an immunometabolic island that
a diagnosis or a severity score cannot see.
